\documentclass[aspectratio=169, 10pt]{beamer}

% --- Theme & Aesthetics ---
\usetheme{default}
\usecolortheme{whale} % Professional blue/grey scheme
\usefonttheme{serif}
\setbeamertemplate{navigation symbols}{}
\setbeamertemplate{itemize items}[circle]

% --- Packages ---
\usepackage{amsmath}
\usepackage{amssymb}
\usepackage{amsthm}
\usepackage{mathtools}
\usepackage{booktabs}
\usepackage{lmodern}
\usepackage{tikz}

% --- Definitions ---
\newcommand{\StateSet}{\mathcal{S}}
\newcommand{\HyphaSet}{\mathcal{H}}
\newcommand{\MushroomSet}{\mathcal{M}}
\newcommand{\PheroField}{\Phi}
\newcommand{\Pressure}{\Psi}
\newcommand{\FlowRate}{Q_f}
\newcommand{\Dimen}{D}
\newcommand{\TaxRate}{\tau}
\newcommand{\RefPrice}{P_{ref}}
\newcommand{\DropThresh}{\delta}
\newcommand{\ProfitTarg}{\rho}

% --- Title ---
\title{Mushroom Brain Model}
\subtitle{Formal specification}
\author{Luis Luarte}
\institute{Agent-Based Systems Modeling}
\date{\today}

\begin{document}
	
	% --- Title Slide ---
	\begin{frame}
		\titlepage
	\end{frame}
	
	% --- Slide 1: Conceptual Narrative ---
	\begin{frame}{Conceptual Foundation: The Fungus and the Market}
		\begin{block}{Resource Allocation by Fungal Networks}
			The model draws its core mechanics from the efficient resource-seeking behavior of fungal mycelia. These networks continuously adapt their topology to optimize nutrient transfer, balancing the costs of network maintenance against the rewards of resource accumulation.
		\end{block}
		
		\vspace{0.3cm}
		
		\textbf{Mapping to Financial Agents:}
		\begin{itemize}
			\item \textbf{Nutrients / Substrate:} Profit Potential.
			\item \textbf{Hyphal Tips ($\mathcal{H}$):} Active trading agents (risk-takers).
			\item \textbf{Mushroom Bodies ($\mathcal{M}$):} Colony sinks; control hyphal-tip exploration.
			\item \textbf{Pheromone Field ($\PheroField$):} A map of successful exploration, guiding agent deployment.
		\end{itemize}
	\end{frame}
	
	% --- Slide 2: Global State Definition ---
	\begin{frame}{System State Space $\StateSet(t)$}
		The simulation evolves over discrete time steps $t \in \mathbb{N}$, defined by the tuple of all active system components.
		
		\begin{equation*}
			\StateSet(t) = \langle W(t), P(t), \PheroField(t), \HyphaSet(t), \MushroomSet(t) \rangle
		\end{equation*}
		
		\vspace{0.3cm}
		\textbf{Key Components:}
		\begin{itemize}
			\item $W(t)$: Global capital pool (System Liquidity).
			\item $P(t)$: Market price signal, $P: \mathbb{N} \to \mathbb{R}^{+}$.
			\item $\PheroField(t)$: The multi-dimensional environment field, $\PheroField: \mathbb{R}^n \to \mathbb{R}^{+}$, representing exploration success.
			\item $\HyphaSet(t)$: The finite set of active Hyphal Agents.
			\item $\MushroomSet(t)$: The finite set of Mushroom Sinks.
		\end{itemize}
	\end{frame}
	
	% --- Slide 3: Hyphal Dynamics I: Efficiency and Flow ---
	\begin{frame}{Hyphal Dynamics I: Efficiency and Flow}
		Hyphal Tips operate in a parameter space $\mathbb{R}^n (n = 2, \text{for first iteration})$. Their trajectory defines their efficiency.
		
		\begin{block}{Path Complexity and Fractal Dimension ($\Dimen$)}
			The efficiency of capital transport is inversely proportional to the complexity of the agent's path history. We define the Fractal Dimension ($\Dimen$) as a metric for path convolution:
			\begin{equation*}
				\Dimen(t) \approx \frac{\log(\text{Total Path Length})}{\log(\text{Net Displacement})}
			\end{equation*}
			where $\Dimen \in [1, 2]$ (1=straight path, 2=area-filling).
		\end{block}
		
		\begin{block}{Flow Rate ($\FlowRate$)}
			The Flow Rate dictates the agent's trading volume and movement velocity:
			\begin{equation*}
				\FlowRate(\Dimen) = 1 - \frac{1}{2^{\Dimen}} \quad \in [0.5, 1.0)
			\end{equation*}
			High complexity (high $\Dimen$) results in higher flow rate (exploit surroundings).
		\end{block}
	\end{frame}
	
	% --- Slide 4: Hyphal Dynamics II: Hydraulic Pressure ---
	\begin{frame}{Hyphal Dynamics II: Hydraulic Pressure ($\Pressure$)}
		$\Pressure$ is the agent's internal state variable, representing urgency and resource status. It drives both movement and resource sharing.
		
		\begin{block}{Internal Pressure Definition}
			$\Pressure$ is a balance between realized gain (weighted by $\beta_1$) and unrealized loss (weighted by $1-\beta_1$, normalized by time constant $\tau$).
			\begin{equation*}
				\Pressure(t) = \beta_{1} \cdot (\text{Realized PnL} \cdot \FlowRate) - (1-\beta_{1}) \cdot \frac{|\text{Unrealized PnL}|}{\tau}
			\end{equation*}
			where $\beta_1 \in [0, 1]$ is the greed factor.
		\end{block}
		
		\vspace{0.3cm}
		\textbf{Behavioral Bifurcation:}
		\begin{itemize}
			\item If $\Pressure \ge \Pressure_{crit}$: \textbf{High Metabolism.} The agent is active, potentially sporulating, and remits capital.
			\item If $\Pressure < \Pressure_{crit}$: \textbf{Stress/Liquidation.} The agent halts movement and trading, entering a death cascade if $\Pressure$ does not recover.
		\end{itemize}
	\end{frame}
	
	% --- Slide 5: Trading Strategy: Spatially Encoded DCA ---
	\begin{frame}{Trading Strategy: Spatially Encoded DCA}
		Each agent's spatial location $x = (\DropThresh, \ProfitTarg)$ dictates its core Dollar-Cost Averaging (DCA) parameters.
		
		\vspace{0.3cm}
		\begin{columns}
			\column{0.5\textwidth}
			\textbf{1. Drop Threshold ($\DropThresh$)}
			The relative percentage price drop required to trigger a new purchase order ($k$).
			\begin{equation*}
				P(t) \le \RefPrice \cdot (1 - \DropThresh \cdot \mathcal{V}_{mult}^{k})
			\end{equation*}
			\textit{Order Size:} Scales with $\FlowRate$.
			
			\column{0.5\textwidth}
			\textbf{2. Profit Target ($\ProfitTarg$)}
			The required gain over the weighted average purchase price ($P_{avg}$) to trigger a sell order.
			\begin{equation*}
				P(t) \ge P_{avg} \cdot (1 + \ProfitTarg)
			\end{equation*}
		\end{columns}
	\end{frame}
	
	% --- Slide 6: Colonial Coordination ---
	\begin{frame}{Colonial Coordination: Resource Sinks and Quorum Sensing}
		\begin{block}{Resource Sinks and Taxation}
			A Mushroom Body $m \in \MushroomSet$ acts as a centralized capital sink. It applies a pressure vacuum $\Pressure_{vac}$ on its lineage.
			
			\begin{itemize}
				\item \textbf{Remittance Flow ($J$):} If $\Pressure_h > \Pressure_{vac}$, the Hypha remits profit:
				\[ J_{h\rightarrow m} = \TaxRate \cdot (\Pressure_h - \Pressure_{vac}) \]
				\item \textbf{Liquidation:} If $\Pressure_h < \Pressure_{vac}$ for $k$ consecutive steps, the agent is dissolved, and capital is recycled to $W$.
			\end{itemize}
		\end{block}
		
		\vspace{0.1cm}
		\begin{block}{Sporulation and Germination}
			The Mushroom sporulates new agents (Spores $\mathcal{S}$) when its mass $M_{mass}$ reaches a threshold. A spore $s$ only germinates (differentiates into a new colony $\MushroomSet'$) if:
			\begin{equation*}
				\PheroField(x_s) > \PheroField_{crit}
			\end{equation*}
			This ensures colonization only occurs in parameter spaces previously proven successful.
		\end{block}
	\end{frame}
	
\end{document}